\section{Classic Protocols to Handle Resource Sharing}
It is stated by~\cite{Buttazzo.2024} that the first four protocols, which are mentioned above, have been developed under the fixed priority assignments but some of them have also been extended under the EDF; whereas, the last one is applicale for both fixed and dynamic priority assignments.
\subsection{Non-Preemptive Protocol (NPP)}
In non-preemptive scheduling, once a task starts executing, it cannot be interrupted until it completes or voluntarily yields the CPU. This simplifies resource access since tasks complete before others interfere.
\begin{itemize}
    \item Benefits: 
    \begin{itemize}
        \item No priority inversion by design.
        \item No need for locks in many simple systems.
    \end{itemize}
    \item Drawbacks:
    \begin{itemize}
        \item High-priority tasks may suffer long blocking times (poor responsiveness)
        \item Not scalable for complex systems with mixed criticality
    \end{itemize}
\end{itemize}
\subsection{Priority Inheritance Protocol (PIP)}
The Priority Inheritance Protocol (PIP) is a classic approach to reduce priority inversion in systems where tasks share mutually exclusive resources (e.g., locks, semaphores). When a low-priority task holds a resource required by a higher-priority task, the low-priority task temporarily inherits the higher priority. This prevents it from being preempted by intermediate-priority tasks, allowing it to finish quickly and release the resource.
\begin{itemize}
    \item Benefits: 
    \begin{itemize}
        \item Reduces unbounded priority inversion
        \item Easy to implement in most RTOS
    \end{itemize}
    \item Drawbacks:
    \begin{itemize}
        \item Can cause deadlocks if resource ordering is not handled carefully
        \item Inheritance must be handled recursively for nested resource requests
    \end{itemize}
\end{itemize}
\subsection{Priority Ceiling Protocol (PCP)}
The Priority Ceiling Protocol (PCP) assigns a ceiling priority to each shared resource — equal to the highest priority of any task that might access it.
A task can only lock a resource if its priority is higher than the ceiling of any currently locked resource by another task.
When a task accesses a resource, it inherits the ceiling priority, blocking any higher-priority tasks that could interfere.
\begin{itemize}
    \item Benefits: 
    \begin{itemize}
        \item Prevents both priority inversion and deadlocks
        \item Ensures bounded blocking time
    \end{itemize}
    \item Drawbacks:
    \begin{itemize}
        \item Requires knowledge of all task/resource relationships at design time
        \item Can be overly restrictive in dynamic systems
    \end{itemize}
\end{itemize}
\subsection{Stack Resource Policy (SRP)}
The Stack Resource Policy (SRP) is used with preemptive scheduling (like EDF or RM) and assumes a single stack for all tasks.
Each resource has a ceiling, and a task can only preempt the current task if its priority is higher than the system ceiling (i.e., the ceiling of all currently locked resources).
\begin{itemize}
    \item Benefits: 
    \begin{itemize}
        \item No runtime deadlocks
        \item Bounded blocking time
    \end{itemize}
    \item Drawbacks:
    \begin{itemize}
        \item Works best with static task sets
        \item More complex to analyze than PIP
    \end{itemize}
\end{itemize}